\section{Appendix C: Testing reductions and proofs of the main results}

The normalized comparison in \cref{obj:thm:normalized-converse} supplies separated treatment effects under nearby observed-sample mixtures. The auxiliary results below translate that comparison into lower bounds for squared-error risk and honest expected interval length. A parametric subexperiment supplies the sampling floors throughout the parameter range, while the normalized priors supply the dimension contribution under their stated regime condition.

\subsection{Parametric sampling floors}

A one-label, complete-arrival Bernoulli subexperiment belongs to \(\mathcal M(d,q)\), the law class in \cref{obj:def:law-class}, for every \(d\ge1\) and \(q\in[1/2,1]\). Its sampling uncertainty supplies the parametric component of both statistical criteria. The underlying two-hypothesis testing comparison is standard \citep[Theorem~2.2]{Tsybakov2009}. The following statement records the resulting floors uniformly over the support bound and arrival floor.

\begin{lemmav}{lem:parametric-floor-infrastructure}[Parametric risk and length floors]
There exists a universal constant \(c_0>0\) such that, for every pair of positive integers \(n,d\) and every \(q\in[1/2,1]\), the minimax squared-error risk \(\mathfrak R^*_{n,d,q}\) defined in \cref{obj:def:point-risk} satisfies
\[
\mathfrak R^*_{n,d,q}\ge \frac{c_0}{n}.
\]
For each noncoverage level \(\alpha\in(0,1/2)\), there exists a constant \(c_{0,\alpha}>0\), depending only on \(\alpha\), such that, for every pair of positive integers \(n,d\) and every \(q\in[1/2,1]\), the minimax worst-law expected honest interval length \(\mathfrak L^*_{n,d,q,\alpha}\) defined in \cref{obj:def:length-risk} satisfies
\[
\mathfrak L^*_{n,d,q,\alpha}\ge \frac{c_{0,\alpha}}{\sqrt n}.
\]
\end{lemmav}

These bounds retain ordinary sampling uncertainty even at complete arrival. For estimation, the floor is universal; for honest intervals, its constant accommodates the fixed coverage requirement. They supply the lower bounds when the missingness-and-dimension contribution lies below the sampling scale and remain available in the regime covered by the normalized converse.

\subsection{From normalized priors to statistical lower bounds}

For each fixed mixture tolerance \(\beta\in(0,1/4)\), choose a positive separation constant \(c_\beta\) satisfying \cref{obj:thm:normalized-converse}. Fix this choice before varying \(n,d,q\); its uniformity is over the parameter triples satisfying that statement, including the condition \(g_{n,d,q}^{\,2}\ge n^{-1}\).

The testing-error comparison applies to the two complete observed-sample mixtures as probability laws \citep[Theorem~2.2]{Tsybakov2009}. The model-specific input is their construction and target separation in \cref{obj:def:paired-prior-handle,obj:thm:normalized-converse}. For squared-error risk, separation is measured in squared units. For honest connected intervals, coverage of separated target values links the same comparison to expected length. The next statement fixes the tolerances for these two uses and records their numerical losses.

\begin{lemmav}{lem:prior-reduction}[Prior reduction and lower bounds]
Let \(c_\beta\) denote the fixed separation constant selected from the normalized converse at mixture tolerance \(\beta\in(0,1/4)\), and set \(\beta_0=1/16\). Write
\[
\ell_n=\log(e+n),
\qquad
g_{n,d,q}=(1-q)\min\left\{1,\frac{d}{n\ell_n}\right\}.
\]
Use the minimax risks and rate in \cref{obj:def:point-risk,obj:def:length-risk,obj:def:rate}.

There exist universal constants \(c>0\) and \(c_0>0\) such that, for every \(\alpha\in(0,1/2)\), there exist constants \(c_\alpha>0\) and \(c_{0,\alpha}>0\) with the following guarantees. Set
\[
\beta_\alpha=\frac{1-2\alpha}{8}.
\]
For every \(n,d,q\) satisfying
\begin{itemize}
\item \textbf{(Sample size.)} \(n\) is an integer with \(n\ge1\).
\item \textbf{(Support bound.)} \(d\) is an integer with \(d\ge1\).
\item \textbf{(Arrival floor.)} \(q\in[1/2,1]\).
\end{itemize}
whenever \(g_{n,d,q}^2\ge n^{-1}\), both bounds
\[
\mathfrak R^*_{n,d,q}
\ge \frac{13}{32}c_{\beta_0}^{\,2}g_{n,d,q}^{\,2},
\qquad
\mathfrak L^*_{n,d,q,\alpha}
\ge \frac54 c_{\beta_\alpha}(1-2\alpha)g_{n,d,q}
\]
hold. Throughout the stated parameter range, the parametric bounds are
\[
\mathfrak R^*_{n,d,q}\ge\frac{c_0}{n},
\qquad
\mathfrak L^*_{n,d,q,\alpha}\ge\frac{c_{0,\alpha}}{\sqrt n}.
\]
The combined guarantees are
\(\mathfrak R^*_{n,d,q}\ge c\,r_{n,d,q}\) and
\(\mathfrak L^*_{n,d,q,\alpha}\ge c_\alpha\sqrt{r_{n,d,q}}\).
\end{lemmav}

The fixed tolerance \(\beta_0\) yields a universal squared-risk constant. By contrast, the coverage-dependent tolerance \(\beta_\alpha\) allocates the mixture-comparison error relative to the positive coverage margin \(1-2\alpha\). For each fixed \(\alpha\in(0,1/2)\), it belongs to the tolerance range required by \cref{obj:thm:normalized-converse}, and the resulting length constant depends only on \(\alpha\).

The displayed coefficients retain the losses associated with observed-mixture proximity and the prior probabilities of target separation. Their role is to make the conversion from the normalized converse to the statistical criteria explicit. Connectedness is relevant to the length criterion because an interval containing separated values spans their distance.

The dimension bounds retain the condition \(g_{n,d,q}^{\,2}\ge n^{-1}\) from \cref{obj:thm:normalized-converse}. The parametric floors in \cref{obj:lem:parametric-floor-infrastructure} supply the sampling contribution throughout the parameter range, including \(q=1\). Together, these inputs give the combined guarantees for the benchmark in \cref{obj:def:rate}, with universal constants for squared risk and constants depending only on fixed noncoverage for expected length.

\subsection{Completion of the main comparisons}

The upper bounds come from \cref{obj:lem:upper-risk,obj:lem:honest-interval-construction}; the normalized-prior reduction and parametric floors in \cref{obj:lem:prior-reduction,obj:lem:parametric-floor-infrastructure} give the matching lower bounds. The formal proofs of the two frontier theorems and the complete-arrival endpoint appear in the following proof appendix, so no second derivation is repeated here.
